\subsection{Koha mesatare e daljes}
Për bredhjen simetrike në nyjet $0,1,\ldots,L$ me kufij thithës, le të jetë $T_i$ numri mesatar i hapave deri në dalje kur nisemi nga $i$. Kushtëzimi ndaj hapit të parë jep
\begin{equation}
 T_i=1+\frac12T_{i-1}+\frac12T_{i+1},
 \qquad T_0=T_L=0.
\end{equation}
Riorganizimi jep diferencën e dytë
\begin{equation}
 T_{i+1}-2T_i+T_{i-1}=-2.
\end{equation}
Zgjidhja kuadratike që plotëson kufijtë është
\begin{equation}
 T_i=i(L-i).
\end{equation}
Në kufirin e difuzionit, koha mesatare $\tau(x)$ plotëson ekuacionin e prapëm
\begin{equation}
 D\tau''(x)=-1,qquad \tau(0)=\tau(a)=0,
\end{equation}
me zgjidhje $\tau(x)=x(a-x)/(2D)$. Kjo lidh drejtpërdrejt ekuacionin diskret me modelin e vazhdueshëm.

Nëse një kufi është reflektues, kushti bëhet derivat zero. Me drift, operatori i prapëm është $v\partial_x+D\partial_{xx}$. Kohët e para të kalimit janë të ndjeshme ndaj bishtave dhe kërkojnë shumë realizime; mesatarja vetëm mund të mos përshkruajë shpërndarjen.

\subsubsection{Spektri i matricës së kalimit}
Për një zinxhir të kthyeshëm, matrica e kalimit është vetë-adjunkte në prodhimin skalar të peshuar
\begin{equation}
 \langle f,g\rangle_\pi=\sum_i\pi_if_ig_i.
\end{equation}
Vlerat vetjake janë reale dhe plotësojnë $1=\lambda_1>\lambda_2\ge\cdots\ge-1$. Komponentët jo stacionarë shuhen si $\lambda_k^n$. Boshllëku spektral $1-\max_{k>1}|\lambda_k|$ kontrollon relaksimin më të ngadaltë.

Për një observabël që projektohet fort mbi modin e dytë, autokorrelacioni mund të sillet si $\rho(n)\simeq\lambda_2^n$. Atëherë
\begin{equation}
 \tau_{\mathrm{int}}\simeq\frac12+\frac{\lambda_2}{1-\lambda_2}
 =\frac{1+\lambda_2}{2(1-\lambda_2)}.
\end{equation}
Kjo tregon lidhjen ndërmjet përzierjes, boshllëkut spektral dhe gabimit të mesatares.

\subsubsection{Vlerësimi praktik i autokorrelacionit}
Vlerësuesi i autokovariancës në vonesën $k$ është
\begin{equation}
 \widehat C(k)=\frac{1}{N-k}\sum_{n=1}^{N-k}
 (A_n-\overline A)(A_{n+k}-\overline A).
\end{equation}
Për $k$ të madh mbeten pak çifte dhe vlerësuesi bëhet i zhurmshëm. Shuma naive deri në $N$ mund të luhatet fort. Përdoret një dritare vetëkonsistente: shuma ndalet kur $k$ është disa herë më i madh se $\widehat\tau_{\mathrm{int}}$, ose kur autokorrelacioni hyn në një rajon të pajtueshëm me zhurmën.

Metoda e blloqeve ndan serinë në blloqe me gjatësi $B$ dhe llogarit mesataren e çdo blloku. Kur $B\gg\tau$, mesataret e blloqeve janë afërsisht të pavarura. Gabimi i vlerësuar duhet të arrijë një pllajë me rritjen e $B$, para se numri i blloqeve të bëhet tepër i vogël.

\subsubsection{Ergodiciteti dhe kurthet}
Stacionariteti nuk mjafton nëse zinxhiri nuk mund të arrijë të gjithë mbështetjen. Një zinxhir i reduktueshëm mund të ketë disa shpërndarje stacionare dhe rezultati varet nga gjendja fillestare. Periodiciteti mund të shkaktojë alternim pa konvergjencë pikë për pikë, edhe kur mesataret kohore sillen më mirë.

Në hapësira me moda të ndara nga barriera, zinxhiri mund të jetë teorikisht ergodik, por praktikisht të mos kalojë barrierën gjatë simulimit. Trajektorja duket e qëndrueshme brenda një mode dhe jep pacaktueshmëri artificialisht të vogël. Nisja e shumë zinxhirëve, propozimet globale, temperimi ose metodat e grupimeve përdoren për të zbuluar dhe zbutur këtë problem.

\subsubsection{Renditja e faqeve dhe zinxhiri Markov}
Në renditjen e faqeve, kalimi ndjek një lidhje me probabilitet $\alpha$ dhe ``teleportohet'' me probabilitet $1-\alpha$. Termi i teleportimit e bën matricën pozitive, kështu që teorema Perron--Frobenius garanton një vektor stacionar unik pozitiv. Parametri $\alpha$ balancon respektimin e strukturës së lidhjeve me shpejtësinë e përzierjes.

Faqet pa dalje duhet të trajtohen: kolona përkatëse zëvendësohet me një shpërndarje probabilitare. Pa këtë hap, matrica nuk është stohastike dhe masa humbet. Ky shembull tregon se pastrimi i të dhënave është pjesë e modelit matematik, jo vetëm detaj programimi.
